\documentclass[10pt,a4paper]{amsart}
\usepackage[margin=19mm]{geometry}
\usepackage{amsmath,amssymb,mathtools}
\usepackage[scr=rsfs,cal=euler]{mathalfa}
\usepackage{xfrac}
\usepackage[unicode,hidelinks,bookmarks=false]{hyperref}
\newcommand{\ie}{i.e.\ }
\newcommand{\cf}{cf.\ }
\newcommand{\resp}{respectively\ }
\newcommand{\Subsection}{\subsection}
\providecommand{\nobreakdash}{\nobreak\hbox{-}}
\setlength{\emergencystretch}{2em}
\multlinegap0pt
\usepackage{amssymb}
\newtheorem{theorem}{Theorem}
\title{On the Real Spectum of the One-Dimensional Dirac Operator with PT-Symmetric Coefficients}
\author{O. A. Veliev}
\date{Source: arXiv:2608.27059v1 (CC BY 4.0)}
\begin{document}
\maketitle

\noindent\emph{Source and licence.} On the Real Spectum of the One-Dimensional Dirac Operator with PT-Symmetric Coefficients. Version arXiv:2608.27059v1 (CC BY 4.0), distributed by arXiv under the Creative Commons Attribution 4.0 International licence, http://creativecommons.org/licenses/by/4.0/, as recorded in the arXiv licence metadata for this exact version and shown on the abstract page https://arxiv.org/abs/2608.27059v1. Full licence notice: \texttt{licenses/arxiv-2608.27059-CC-BY-4.0.txt}.\par\smallskip

\noindent\textit{Editorial scope.} Counted unit: the article's theorem theorem (source0.tex lines 355--358). The article's complete original proof of it is delivered with it (source0.tex lines 360--414, one \texttt{proof} environment of 55 lines). No mathematics is written, repaired or re-derived: the statement and the proof are the article's own lines, in the article's own notation, and no retained line is substituted or re-flowed.  No further source-local context is needed: every cross-reference of every retained line resolves inside the two delivered spans.  Nothing was omitted from any retained span, so the package carries no omission record.  No retained line contains a citation, so the wrapper carries no bibliography and no bibliography entry.  No retained line contains a figure environment, an image-inclusion command or drawing code, so no image is delivered and the package declares no asset dependency.  Editorial formatting only, in addition: an AMS \texttt{amsart} preamble that restates the article's own declarations and macros used by the retained lines, namely the environments \texttt{theorem} on separate counters, together with the standard packages those lines need.  The retained environments print in this document as Theorem 1 in order of appearance, whereas the article prints the same items with the numbering of its own full document.  The complete native source of the article as distributed in the arXiv e-print for this exact version is bundled unchanged as \texttt{sources/208680/source0.tex}.
\begin{theorem}
The expression $l(n,m)$ and the operator $L(n,m)$ are PT-symmetric if and only
if $p_{k,i,j}$ are PT-symmetric functions for all $k,i$, $j.$
\end{theorem}

\begin{proof}
\bigskip Let $p_{k,i,j}$ be PT-symmetric functions for all $k,i$, $j$. The
operator $L(n,m)$ can be written in the form
\[
L(n,m)\mathbf{y}=\sum_{k=0}^{n}P_{k}D^{n-k}\mathbf{y},
\]
where $D\mathbf{y}=-i\mathbf{y}^{^{\prime}}.$ It is enough to prove that
\begin{equation}
\left(  P_{k}D^{n-k}\right)  PT\mathbf{y}(x)=PT\left(  P_{k}D^{n-k}%
\mathbf{y}(x)\right)  \tag{10}%
\end{equation}
for each $k=0,1,...,n$ and for all $\mathbf{y}\in D(L(n,m)).$ To this end, let
us calculate the left- and right-hand sides of (10) separately. Since
$PT\mathbf{y}(x)=\overline{\mathbf{y}(-x)}$ , by the chain rule, we have
$\left(  PT\mathbf{y}(x)\right)  ^{^{\prime}}$ $=\left(  \overline
{\mathbf{y}(-x)}\right)  ^{^{\prime}}$ $=-\overline{\mathbf{y}^{\prime}(-x)}$
and $D\overline{\mathbf{y}(-x)}=i\left(  \overline{\mathbf{y}^{\prime}%
(-x)}\right)  .$ Hence,
\[
\text{ }\left(  P_{k}(x)D^{n-k}\right)  PT\mathbf{y}(x)=i^{n-k}P_{k}%
(x)\overline{\mathbf{y}^{(n-k)}(-x)}.
\]
On the other hand, for the right-hand side of (10), we have%
\begin{align*}
PT\left(  P_{k}(x)D^{n-k}\mathbf{y}(x)\right)   &  =PT\left(  P_{k}%
(x)(-i)^{n-k}\mathbf{y}^{(n-k)}(x)\right)  =\\
P(\overline{P_{k}(x)}(i)^{n-k}\overline{\mathbf{y}^{(n-k)}(x)})  &
=(\overline{P_{k}(-x)}(i)^{n-k}\overline{\mathbf{y}^{(n-k)}(-x)}).
\end{align*}
Therefore, (10) holds due to the equality $\overline{P_{k}(-x)}=P_{k}(x).$
Thus, the sufficiency is proved.

Now suppose that $\overline{P_{k}(-x)}-P_{k}(x)\neq O$ for $k\in\left\{
k_{\,1},k_{2},...,k_{s}\right\}  ,$ but $L(n,m)$ is a PT-symmetric operator,
where $0\leq k_{\,1}<k_{2}<...<k_{s}\leq n.$ Then, by Definition 1, we have
\[
PTL(n,m)\mathbf{y}(x)-L(n,m)PT\mathbf{y}(x)\mathbf{=}\sum_{j=1}^{s}%
(i)^{n-k_{j}}\left(  (\overline{P_{k_{j}}(-x)}-P_{k_{j}}(x))\overline
{\mathbf{y}^{n-k_{j}}(-x)}\right)  .
\]
Therefore, the equality
\begin{equation}
PTL(n,m)\mathbf{y}(x)=L(n,m)PT\mathbf{y}(x) \tag{11}%
\end{equation}
is satisfied only for those functions $\mathbf{y}(x)$ for which $\overline
{\mathbf{y}(-x)}$ is a solution of the equation
\[
\sum_{j=1}^{s}(i)^{n-k_{j}}\left(  (\overline{P_{k_{j}}(-x)}-P_{k_{j}%
}(x))\overline{\mathbf{y}^{n-k_{j}}(-x)}\right)  =0.
\]
Since the set of solutions of this equation is a finite dimensional subset $E$
of $D(L(n,m))$, (11) is not satisfied for any $\mathbf{y}\in\left(
D(L(n,m))\backslash E\right)  .$ This means that $L(n,m)$ is not a
PT-symmetric operator. The theorem is proved.
\end{proof}

\end{document}
